% !TeX root = meccanica_analitica.tex
\documentclass{article}

\usepackage[T1]{fontenc}
\usepackage[utf8]{inputenc}
\usepackage[a4paper, total={15.3 cm, 21.3 cm}]{geometry}
\usepackage{amsmath}
\usepackage{amssymb}
\usepackage{amsthm}
\usepackage{gensymb}
\usepackage{booktabs}
\usepackage{hyperref}
\usepackage{caption}
\usepackage{float}
\usepackage{graphicx}
\usepackage{mathtools}
\usepackage{subfig}
\usepackage{titlesec}
\usepackage{titletoc}
\usepackage{physics}
\usepackage{siunitx}
\usepackage[dvipsnames]{xcolor}

\usepackage{longtable}
\usepackage{tabularx}
\usepackage{calc}
\usepackage{array}
\usepackage{subfiles}
\usepackage{etoolbox}
\usepackage{xparse}

\newtheorem{theorem}{Theorem}[section]
\newtheorem{corollary}{Corollary}[theorem]
\newtheorem{lemma}{Lemma}[theorem]
\newtheorem{definition}{Definition}[section]
\newtheorem{example}{Example}[section]
\newtheorem{proposition}{Proposition}[section]

\hypersetup{colorlinks=true,linkcolor=black}
\renewcommand\thesection{\arabic{section}}
\titlecontents{chapter}[1.05em]{\bigskip}
{\contentslabel[\MakeUppercase{\romannumeral\thecontentslabel}]{1em}\enspace\textsc}
{\hspace*{-1em}\textsc}
{\hfill\contentspage}
\titlecontents{section}[1.6em]{\smallskip}
{\thecontentslabel.\enspace}
{}
{\titlerule*[1pc]{.}\contentspage}
\setcounter{tocdepth}{2}


\begin{document}

    \pagenumbering{roman}
    \thispagestyle{empty}

    % First page with uni logo and title
    \begin{center}

        \includegraphics[width=1.\linewidth]{unimi_logo.jpg}
        \centering
        \vspace{3cm}

        \uppercase{\Large MECCANICA ANALITICA: \\ APPUNTI DEL CORSO \par}
        \vspace{3cm}

        \Large Lorenzo Liuzzo\par
        \vspace{1.5cm}

        \Large \today

    \end{center}

    \clearpage

    % Table of contents
    \tableofcontents

    \clearpage
    \pagenumbering{arabic}

    \section{Lagrangiana}

        \subsection{Problema ad un corpo}
        
            \begin{theorem}(della forza viva)
                \label{thm:Forza viva}

                Sia $T$ la forza viva (o energia cinetica) come $T = \frac{1}{2}m \dot{x} \cdot \dot{x}$.
                Allora lungo ogni soluzione $x = x(t)$ dell'equazione di Newton $m \ddot{x} = F(x)$, si ha 
                \[ \dot{T} = F(x) \cdot \dot{x} \] o equivalentemente \[ T(t_1) - T(t_0) = \int_{t_0}^{t_1} F(x) \cdot \dot{x} \, dt \]

                dove $F(x) \cdot \dot{x}$ è la potenza della forza e $\int_{t_0}^{t_1} F(x) \cdot \dot{x}$ è il lavoro svolto dalla forza. 

            \end{theorem}
            \begin{proof}

                Moltiplicando l'equazione di Newton per $\dot{x}$ e applicando la regola di Leibniz per la derivata di un prodotto, si ottiene
                \[ (m \ddot{x}) \cdot \dot{x} = \frac{d}{dt}(\frac{1}{2} m \dot{x} \cdot \dot{x}) = \dot{T} \]

            \end{proof} 

            Nel caso di un campo di forze posizionali $F = F(x)$ l'integrale a secondo membro dipende dal movimento $x(t)$ nell'intervallo $(t_0, t_1)$ 
            attraverso la corrispondente traiettoia $\gamma$. Si ha dunque un integrale curvilineo della forma differenziale: 
            \[ T(t_1) - T(t_0) = \int_{\gamma} F(x) \cdot dx \] \\
            Nel caso in cui la forza sia conservativa, cioè in cui la forza ammetta potenziale, ossia una funzione scalare $V = V(x)$ tale che $F = - grad(V)$, 
            chiamiamo $V$ energia potenziale (o funzione delle forze). Questa condizione può essere riscritta dicendo che la forma differenziale del lavoro 
            è esatta, cioè è il differenziale di una funzione, in particolare: $F(x) \cdot dx = -dV(x)$. \\

            \begin{theorem}(dell'energia)
                \label{thm:Conservazione dell'energia}

                Per un punto soggetto ad un campo di forze posizionali conservativo $F = F(x)$, lungo ogni soluzione $x = x(t)$ dell'equazione di Newton $m \ddot{x} = F$, 
                si ha \[ E = T + V \quad \dot{E} = 0 \]
                
            \end{theorem}
            \begin{proof}

                \[ F \cdot \dot{x} = - \frac{\partial V}{\partial x} \cdot \frac{dx}{dt} = - \frac{dV}{dt} = - \dot{V} \quad \Longrightarrow \quad \dot{T} = -\dot{V} \]
                oppure, sfruttando la definizione integrale di lavoro (\ref{thm:Forza viva}) e il fatto che l'integrale curvilineo di una forma differenziale esatta
                lungo una curva orientata dipenda solo dai suoi estremi (cioè che $\oint F(x) \cdot dx = 0$), si ha
                \[ \int_{\gamma}F(x) \cdot dx = - \int_{\gamma} dV = V(A) - V(B) \]

            \end{proof}
            
            \begin{definition}(costante del moto)
                \label{def:costante del moto}

                Una variabile dinamica che assume un medesimo valore per ogni punto del moto corrispondente ad una soluzione dell'equazione di Newton è detta \boldmath{costante del moto}.
            
            \end{definition}
            Dunque, la funzione $E(x, \dot{x}) = T(\dot{x}) + V(x)$ è una costante del moto. Pertanto, fissati i dati iniziali $x_0$ e $\dot{x}_0$, e quindi anche $E_0$,
            il teorema di conservazione dell'energia totale (\ref{thm:Conservazione dell'energia}) va inteso nella forma \[ T - V = E_0 = E(x, \dot{x})\] \\
            Introducendo la quantitità di moto $p = m \dot{x}$, è possibile riscrivere dell'equazione di Newton come 
            \[ \dot{p} = F \quad \longleftrightarrow \quad p(t_1) - p(t_0) = \int_{t_0}^{t_1} F dt \] dove l'integrale della forza nel tempo viene detto impulso della forza. 

            \begin{theorem}(del momento angolare)
                \label{thm:momento angolare}

                Per un punto materiale soggetto ad una generica forza $F = F(x, \dot{x}, t)$, lungo ogni soluzione $x = x(t)$ dell'equazione di Newton $m \ddot{x} = F$, si ha
                \[ \dot{L} = M \] dove $L = x \times p$ è il momento angolare e $M = x \times F $ è il momento angolare della forza.
                
            \end{theorem}
            \begin{proof}
                
                Moltiplicando vettorialmente l'equazione di Newton per $x$ e applicando la regola di Leibniz per la derivata di un prodotto, si ottiene
                \[ M = x \times F = x \times m \ddot{x} = x \times \frac{d}{dt}(m\dot{x}) = \frac{d}{dt}(x \times p) - \frac{dx}{dt} \times p = \frac{d}{dt}(x \times p) = \dot{L} \]
           
            \end{proof}

            \begin{corollary}(conservazione del momento angolare)
                \label{cor:conservazione del momento angolare}

                Per un punto materiale soggetto ad un campo di forze centrali $F = F(x)$, lungo ogni soluzione $x = x(t)$ dell'equazione di Newton $m \ddot{x} = F$, si ha
                \[ \dot{L} = 0 \] ossia che il momento angolare $L$ è una costante del moto.

            \end{corollary}
            \begin{proof}

                Si applica il teorema del momento angolare (\ref{thm:momento angolare}). Poichè $F$ è un campo di forze centrali, si ha che la forza è parallela al raggio vettore,
                perciò $ M =  x \times F = 0 \Longrightarrow \dot{L} = 0 $.

            \end{proof}

            \begin{corollary}(campi centrali e moti bidimensionali)
                \label{cor:campi centrale e moti bidimensionali}

                Per un punto materiale soggetto ad un campo di forze centrali $F = F(x)$, per ogni soluzione $x = x(t)$ dell'equazione di Newton $m \ddot{x} = F$, 
                la traiettoia $x = x(t)$ giace in un piano passante per il centro di forza e ortogonale al vettore $L$ momento angolare, determinato dalle 
                condizioni iniziali $x_0$ e $\dot{x}_0$ che definiscono $L_0 = x_0 \times m \dot{x_0}$.

            \end{corollary}
            \begin{proof}
                
                Per le proprietà del prodotto vettoriale, si ha che $L = x \times p$ è ortogonale a $x$. Ma per la conservazione del momento angolare (\ref{cor:conservazione del momento angolare}),
                si ha che $L$ è una costante del moto. Dunque, per ogni tempo $t$, il vettore $x(t)$, dovendo essere ortogonale ad un vettore costante, giace in un piano ortogonale a quel vettore.

            \end{proof}

            Prendendo l'equazione della quantità di moto $\dot{p} = F$ e proiettandola su un asse qualsiasi, ad esempio $x$, si ottiene $\dot{p}_x = F_x$ e 
            quindi se $F_x = 0 \Longrightarrow \dot{p}_x = 0$ e quindi $p_x = cost$. \\
            Ma per forze posizionali conservative, $F_x = 0$ corrisponde a dire che $V$ è invariante per traslazione lungo l'asse $x$, ossia 
            \[ \frac{\partial V}{\partial x} = 0 \quad \textrm{ovvero che} \quad V(x + h, \dots) = V(x, \dots) \quad \forall h \]
            \begin{proposition}
                \label{prop:simmetrie e conservazioni}

                Se l'energia potenziale $V$ è invariante per traslazioni lungo un asse, allora la componente della quantità di moto $p$ lungo quell'asse è una costante del moto.

            \end{proposition}
            Analogamente per il momento angolare, se ad esempio $M_x = x \cdot \frac{\partial V}{\partial y} - y \cdot \frac{\partial V}{\partial x} = \frac{\partial V}{\partial \varphi} = 0 $, 
            allora $L_x = cost$.
            \begin{proposition}
                \label{prop:simmetrie e conservazioni 2}
                
                Se l'energia potenziale $V$ è invariante per rotazioni rispetto ad un asse, allora la componente del momento angolare $L$ su quell'asse è una costante del moto.
            \end{proposition}
            Inoltre, nei casi in cui la forza sia indipendente dal tempo, cioè $\frac{\partial V}{\partial t} = 0$, si dice che l'energia potenziale è invariante per traslazioni temporali.
            \begin{proposition}
                \label{prop:simmetrie e conservazioni 3}
                
                Se l'energia potenziale $V$ è invariante per traslazioni temporali, allora l'energia $E = T + V$ è una costante del moto.
            \end{proposition} 


        \subsection{Problema a N corpi}
            
            In un sistema di N punti materiali, le N equazioni di Newton del secondo ordine \[ m_k \cdot \ddot{x}_k = F_k \quad (k = 1, 2, \dots, N) \] si possono scrivere come
            un sistema di 2N equazioni di primo ordine \[ \dot{x}_k = \frac{p_k}{m_k} \quad \quad \dot{p}_k = F_k \quad (k = 1, 2, \dots, N) \]
            Si introducono la quantitità di moto totale del sistema, il momento angolare totale del sistema, la risultante delle forze e il momento risultante delle forze
            rispettivamente \[ p = \sum_k p_k = \sum_k m_k \dot{x}_k \] \[ L = \sum_k L_k = \sum_k x_k \times p_k \] \[ R = \sum_k F_k \] \[ M =  \sum_k M_k = \sum_k x_k \times F_k \]
            Le ipotesi sulle forze sono:
            \begin{itemize} 
                \item ogni forza può essere decomposta in una componente dovuta agli altri punti del sistema e una dovuta ad agenti esterni, cioè $F_k = F_k^{(int)} + F_k^{(est)}$.
                \item ogni forza interna può essere decomposta nella somma di forze dovute singolarmente ad ogni altro punto del sistema, cioè $F_k^{(int)} = \sum_{j \neq k} F_{kj}$.
                \item ogni forza rispetti il principio di azione e reazione, cioè $F_{kj} = - F_{jk}$.
                \item ogni forza sia centrale, cioè $F_{kj} = f_{kj}(r_{kj}) \cdot \frac{x_k - x_j}{r_{kj}}$, dove $f_{kj} = f_{jk}$ e $r_{kj} = \Vert x_k - x_j \Vert$.
            \end{itemize}
            Le forze che soddisfano queste ipotesi sono anche dette forze classiche.

            \begin{theorem}(equazioni cardinali della dinamica)
                \label{thm:equazioni cardinali}

                Lungo le soluzioni del sistema di equazioni di Newton per un sistema a N punti materiali soggetto a forze di tipo classico si hanno le seguenti relazioni,
                dette rispettivamente la prima e la seconda equazione cardinale: \[ \dot{p} = R^{(est)} \quad \quad \dot{L} = M^{(est)} \]
                dove $R^{(est)}$ e $M^{(est)}$ sono la risultante e il momento risultante delle forze esterne.

            \end{theorem}
            \begin{corollary}(condizioni necessarie per l'equilibrio di un sistema)
                \label{cor:equilibrio di un sistema}

                Un sistema a N punti materiali soggetto a forze di tipo classico è in equilibrio se e solo se \[ R^{(est)} = 0 \quad \text{e} \quad M^{(est)} = 0 \]

            \end{corollary}

            \begin{theorem}(potenziale delle forze interne)
                \label{thm:potenziale delle forze interne}

                Per forze interne di tipo classico, il lavoro elementare delle forze interne è un differenziale esatto e si ha 
                \[ F^{(int)} \cdot dx = \sum_k F_k^{(int)} \cdot dx_k = -d[\frac{1}{2} \sum_{k, j}^{k \neq j} V_{kj}(r_{kj})] = -dV^{(int)} \]

            \end{theorem}
            \begin{proof}

                Consideriamo una coppia di punti materiali $k$ e $j$. Poichè su essi agiscono soltanto forze di tipo classico ed essendo valido il principio di azione e reazione 
                $F_{kj} = - F_{jk}$, possiamo prendere $r_{kj} = x_k - x_j$ e scrivere \[ F_{kj} \cdot dx_k + F_{jk} \cdot dx_j = F_{kj} \cdot dr_{kj} \]
                ossia ci siamo ridotti al problema di un solo corpo soggetto ad una forza centrale con simmetria sferica.

            \end{proof}

            \begin{theorem}(problema dei due corpi)
                \label{thm:problema dei due corpi}

                Per un sistema di due corpi soggetti soltanto a forze interne, $R^{(est)} = 0$, soddisfacenti il principio di azione e reazione, si ha 
                \[ \ddot{x}_{CM} = 0 \quad \quad \mu \cdot \ddot{r} = F(r) \] dove $\mu$ è la massa ridotta $\mu = \frac{m_1 \cdot m_2}{m_1 + m_2}$ e $r$ è la distanza tra i due corpi.

            \end{theorem}
            \begin{proof}
                
                Dalle equazioni di Newton per i due corpi, sommando membro a membro si ha che \[ m_1 \cdot \ddot{x}_1 + m_2 \cdot \ddot{x}_2 = F_{12} + F_{21} = 0 \] da cui si deduce che 
                \[ \ddot{x}_{CM} = \frac{m_1 \cdot \ddot{x}_1 + m_2 \cdot \ddot{x}_2}{m_1 + m_2} = 0 \]
                Invece, sottraendo membro a membro le due equazioni divise rispettivamente per le masse, si ha \[ \ddot{x}_1 - \ddot{x}_2 = \frac{F_{12}}{m_1} - \frac{F_{21}}{m_2} = 2 * \frac{F(r)}{\mu} \]
                da cui deduciamo \[ \mu \cdot \ddot{r} = F(r) \]

            \end{proof}

            
        \subsection{Spazio ambiente e coordinate generalizzate}
            
            Spesso definiamo globalmente una superficie mediante un'equazione implicita $F(x, y, z) = 0$. Tale rapprestazione impone un vincolo sulle coordinate. 
            \begin{definition}(carta locale e formula di immersione)

                Si ha una rappresentata parametrica locale o una carta locale quando è assegnata una funzione \[ F: U \rightarrow \mathbb{R}^n, \quad U \subseteq \mathbb{R}^{n - 1} \quad \text{aperto} \quad 
                \text{t.c.} \quad x = F(q), \quad q = (q_1, \dots, q_(n - 1)) \] Si dice 'formula di immersione' la funzione $x = x(q)$ che porta dalla carta locale alla superficie nello spazio ambiente.
                
            \end{definition}

            \begin{definition}(linee e vettori coordinati)

                Data una carta $U$ e una superficie $M$, si definiscono linee coordinate quelle linee di $M$ lungo le quali una sola delle coordinate rimane costante, corrispondenti alle linee $q_i = cost$ sulla carta.
                Preso un punto $P \in M$ di coordinate $q = (q_1, \dots, q_n)$ nella carta $x = x(q)$, sono detti vettori coordinati i vettori \[ \frac{\partial x}{\partial q_i} \quad (i = 1, \dots, n) \]
                Tali vettori, spiccati dal punto $P$ e tangenti alle linee coordinate $q_1, \dots, q_n$, sottengono il piano tangente $T_P M$ alla varietà $M$ in $P$, dunque debbono essere linearmenti indipendenti.
                Introducendo la matrice \[ g_{ik} = \frac{\partial x}{\partial q_i} \cdot \frac{\partial x}{\partial q_k} \] questa condizione di lineare indipendenza si traduce in \[ det(g_{ik}) \neq 0 \]
                Per assegnare un movimento $x = x(t)$ sulla superficie $M$, basterà comporlo con la funzione $q = q(t)$ sulla carta locale, in modo tale da ottenere 
                \[ x = x(q(t)) \quad \quad v = \frac{dx}{dt} = \frac{\partial x}{\partial q} \dot{q} = \sum_i \frac{\partial x}{\partial q_i} \dot{q}_i \] 

            \end{definition}

            \begin{theorem}(formula del binomio lagrangiano)
                \label{thm:binomio lagrangiano}

                Vale \[ \ddot{x} \cdot \frac{\partial x}{\partial q_i} = \frac{d}{dt} \frac{\partial \frac{1}{2} \dot{x}^2}{\partial \dot{q}_i} - \frac{\partial \frac{1}{2} \dot{x}^2}{\partial q_i} \quad (i = 1, \dots, n) \]
                che in termini di energia cinetica $T$ diventa \[ m\ddot{x} \cdot \frac{\partial x}{\partial q_i} = \frac{d}{dt} \frac{\partial T}{\partial \dot{q}_i} - \frac{\partial T}{\partial q_i} \quad (i = 1, \dots, n) \] 
                
            \end{theorem}
            \begin{proof}
                
            \end{proof}

            \begin{lemma}
                
                Valgono le relazioni \[ \frac{\partial v}{\partial \dot{q}_i} = \frac{\partial x}{\partial q_i} \quad \quad
                                        \frac{\partial v}{\partial q_i} = \frac{d}{dt} \frac{\partial x}{\partial q_i} \quad \quad (i = 1, \dots, n) \] 

            \end{lemma}
            \begin{proof}
                
                La prima si ricava dal fatto che $v = \frac{\partial x}{\partial q_i} \dot{q}_i$ mentre la seconda segue usando la regola per la derivata del prodotto 
                \[ \frac{d}{dt} \frac{\partial x}{\partial q_i} = \sum_k \frac{\partial}{\partial q_k} (\frac{\partial x}{\partial q_i}) \dot{q}_k 
                                                                = \sum_k \frac{\partial^2 x}{\partial q_k \partial q_i} \dot{q}_k 
                                                                = \frac{\partial}{\partial q_i} (\sum_k \frac{\partial x}{\partial q_k} \dot{q}_k) = \frac{\partial v}{\partial q_i} \]

            \end{proof}
            Pertanto, fissata una carta con formula di immersione $x = x(q)$, si esprime l'energia cinetica in funzione delle $q$ e delle $\dot{q}$ come 
            \[ T = \frac{1}{2} \sum_{i, k} a_{ik}(q) \dot{q}_i \dot{q}_k \quad \quad a_{ik} = \frac{m}{2}g_{ik} = \frac{m}{2} \frac{\partial x}{\partial q_i} \frac{\partial x}{\partial q_k} \]


        \subsection{Equazioni di Lagrange}

            \begin{definition}(vincolo)
                        
            \end{definition}

            \begin{definition}(reazione vincolare)

            \end{definition}

            \begin{definition}(Lagrangiana)
                \label{def:Lagrangiana}

                Si consideri un punto materiale su una superficie rappresentata in una carta locale da una formula di immersione $x = x(q)$ con $q = (q_1, \dots, q_n)$ e si ammetta che valga 
                l'equazione di Newton-d'Alembert \[ m \ddot{x} = F + F^{(v)} \]
                Moltiplicando scalarmente l'equazione di Newton per i vettori coordinati $\frac{\partial x}{\partial q_i} \quad (i = 1, \dots, n) $, poichè essi giacciono nel piano tangente alla superficie 
                ed essendo $F^{(v)}$ ortogonale alla superficie per definizione di reazione vincolare si ha che \[ F^{(v)} \cdot \frac{\partial x}{\partial q_i} = 0 \]
                e dunque \[ m \ddot{x} \cdot \frac{\partial x}{\partial q_i} = F \cdot \frac{\partial x}{\partial q_i} \]
                Nel caso di forze derivanti da potenziale, per cui $ F = - grad(V) $, si ottiene che \[ F \cdot \frac{\partial x}{\partial q_i} = - \frac{\partial V}{\partial q_i} \quad \quad (i = 1, \dots, n) \]
                Rielaborando il primo membro con la formula del binomio lagrangiano \ref{thm:binomio lagrangiano} si ottiene 
                \[ \frac{d}{dt} \frac{\partial T}{\partial \dot{q}_i} - \frac{\partial T}{\partial q_i} = - \frac{\partial V}{\partial q_i} \] 
                Si definisce dunque la funzione lagrangiana $L = L(q, \dot{q})$ come la differenza tra l'energia cinetica e l'energia potenziale del sistema, ossia \[ L = T - V \]
                Si osserva come $L: TM \rightarrow \mathbb{R}$ sia definita intrinsecamente nello spazio degli stati $TM$, il fibrato tangente dello spazio $M$ detto spazio delle configurazioni. 
                
            \end{definition}

            \begin{theorem}(equazioni di Eulero-Lagrange)
                \label{thm:equazioni di Eulero-Lagrange}

                Si consideri un punto materiale la cui posizione sia rappresentata in una carta locale $x = x(q)$ da coordinate $q = (q_1, \dots, q_n)$, allora
                i movimenti $q = q(t)$ sulla carta corrispondenti alle soluzioni dell'equazione di Newton \[ m \ddot{x} = - grad(V) + F^{(v)} \] sono soluzioni delle
                equazioni di Eulero-Lagrange nella carta delle coordinate $q$ \[ \frac{d}{dt} \frac{\partial L}{\partial \dot{q}_i} - \frac{\partial L}{\partial q_i} = 0 \quad \quad (i = 1, \dots, n) \]
                
            \end{theorem}
            \begin{proof}
                
                Proiettando le equazioni di Newton sulle linee coordinate si ha \[ \frac{d}{dt} \frac{\partial T}{\partial \dot{q}_i} - \frac{\partial T}{\partial q_i} = - \frac{\partial V}{\partial q_i} \quad \quad (i = 1, \dots, n) \]
                Definita la lagrangiana $L = L(q, \dot{q})$, poichè $V = V(q)$, si osserva che $\frac{\partial T}{\partial \dot{q}_i} = \frac{\partial L}{\partial \dot{q}_i}$ e quindi si ha la tesi.

            \end{proof}

            \begin{definition}(forza generalizzata)
                
                Nel caso generale in cui la forza attiva dipenda anche dalla velocità o dal tempo, cioè $F = F(q, \dot{q}, t)$, si introduce la forza generalizzata $Q$ come
                \[ Q_i = F \cdot \frac{\partial x}{\partial q_i} \quad \quad (i = 1, \dots, n) \] e le equazioni di Eulero-Lagrange prendono la forma di 
                \[ \frac{d}{dt} \frac{\partial T}{\partial \dot{q}_i} - \frac{\partial T}{\partial q_i} = Q_i \quad \quad (i = 1, \dots, n) \]
                Se consideriamo il lavoro elementare $\delta W$ della forza $F$ si ha \[ \delta W = F \cdot dx = \sum_i Q_i dq_i = Q dq \]
                
            \end{definition}

            \begin{theorem}
                
                Per un sistema a N punti con vincoli perfetti, in una carta locale con coordinate $q = (q_1, \dots, q_n)$, i movimenti $q = q(t)$ sulla carta corrispondenti
                alle soluzioni del sistema di equazioni di Newton-d'Alembert sono soluzioni delle equazioni di Eulero-Lagrange \ref{thm:equazioni di Eulero-Lagrange}
                La condizione geometrica di perfezione del vincolo è equivalente alla condizione che sia nullo il lavoro totale delle reazioni vincolari per ogni spostamento compatibile.
                Infatti, gli spostamenti infinitesimi compatibili con i vincoli si esprimono nella forma \[ dX = \frac{\partial X}{\partial q} dq \] da cui si ottiene l'espressione 
                del lavoro elementare totale delle reazioni vincolari \[ \delta W = \sum_i F_i^{(v)} \cdot dx_i = F^{(v)} \cdot dX = F^{(v)} \cdot \frac{\partial X}{\partial q} dq \]
            \end{theorem}

            \begin{theorem}(dell'energia generalizzata)
                \label{thm:di Jacobi}

                Per un sistema lagrangiano con lagrangiana $L = L(q, \dot{q})$ in una carta locale assegnata, per ogni movimento $q = q(t)$ soddisfacente le equazioni di Eulero-Lagrange \ref{thm:equazioni di Eulero-Lagrange}
                si ha \[ \dot{\mathcal{E}} = - \frac{\partial L}{\partial t} \] dove $\dot{\mathcal{E}} = \dot{\mathcal{E}}(q, \dot{q}, t)$ è l'energia generalizzata del sistema definita come 
                \[ \dot{\mathcal{E}} \coloneqq p\dot{q} - L \quad \quad p \coloneqq \frac{\partial L}{\partial \dot{q}} \]
                In particolare, se la lagrangiana è invariante per traslazioni temporali ossia non dipende la tempo, si ha che l'energia generalizzata è una costante del moto.  

            \end{theorem}

            

    \clearpage
    \section{Hamiltoniana}

    \clearpage
    \section{Principi variazionali}



\end{document}